% Generated draft: requires mathematical and source review.
\documentclass[11pt]{article}
\usepackage[margin=25mm]{geometry}
\usepackage{fontspec}
\setmainfont{Latin Modern Roman}
\usepackage{amsmath,amssymb,mathtools,unicode-math}
\setmathfont{Latin Modern Math}
\usepackage{enumitem,needspace,xurl,hyperref}
\setlength{\parindent}{0pt}
\setlength{\parskip}{5pt}
\setlength{\emergencystretch}{3em}
\newcommand{\N}{\mathbb{N}}
\newcommand{\Z}{\mathbb{Z}}
\newcommand{\Q}{\mathbb{Q}}
\newcommand{\R}{\mathbb{R}}
\newcommand{\C}{\mathbb{C}}
\newcommand{\F}{\mathbb{F}}
\newcommand{\A}{\mathbb{A}}
\newcommand{\PP}{\mathbb{P}}
\newcommand{\E}{\mathbb{E}}
\newcommand{\eps}{\varepsilon}
\newcommand{\dd}{\,\mathrm{d}}
\begin{document}
\section*{UnsolvedMath ID: 3100047}
\textbf{Source problem code:} catalogue entry
\subsection{Definitions and mathematical statement}
Let $\mathcal P(G)$ denote the finite nonempty set of successful pressing sequences of a graph $G$, with the pressing convention and permitted graph class still to be specified. A graph is uniquely pressable if $|\mathcal P(G)|=1$. Whenever the relation ``$x$ occurs before $y$'' is defined for every sequence in $\mathcal P(G)$, put
\[
p_G(x,y)=\frac{|\{\sigma\in\mathcal P(G):x\text{ occurs before }y\text{ in }\sigma\}|}{|\mathcal P(G)|}.
\]
The conjecture asks whether every graph $G$ that is not uniquely pressable has two distinct vertices $x,y$ such that
\[
\frac13\leq p_G(x,y)\leq\frac23.
\]
The random sequence is uniform on $\mathcal P(G)$, and the endpoints are included.

For context, one usual pressing model starts with a finite simple graph whose vertices are black or white. Pressing a black vertex $v$ toggles adjacency between every pair of distinct vertices in its closed neighborhood and toggles the colors of all vertices in that neighborhood. A successful sequence consists of legal presses and ends with an edgeless graph having all vertices white. This convention is not specified in the supplied item: in particular a successful sequence in this model need not contain every vertex. Thus the displayed probability is a precise formulation only after the intended sequence family and the treatment of omitted vertices have been fixed. The item also does not say whether graphs without successful sequences are excluded; uniform sampling requires a nonempty family. These issues are unresolved in the supplied source.
\subsection{Short English statement}
For a graph with more than one successful pressing sequence, must some pair of vertices appear in either relative order with probability at least $1/3$ under the uniform distribution on successful sequences? The pressing model, the eligibility of graphs, and the treatment of vertices omitted from a sequence are not specified in the supplied item.
\subsection{Sources}
\begin{itemize}
\item Cooper, \emph{Combinatorial Problems I Like}, October 2020 snapshot, ``Matroids, Lattices, and Posets,'' source-order bullet 47: \url{https://people.math.sc.edu/cooper/combprob.html}; accessed 2026-07-29, HTML extraction. The item refers to ``See this,'' but the supplied record includes no destination for that link.
\item Local registry \texttt{data/unsolvedmath/problems.json}, permanent UnsolvedMath ID $3100047$, source code catalogue entry, \emph{Is the 1/3-2/3 Conjecture for Pressing Sequences true}. Dataset attribution: Ulam AI and the UnsolvedMath Contributors.
\item The supplied research summary connects the question to pressing-sequence research of Cooper and collaborators and to the poset $1/3$--$2/3$ conjecture; it supplies no specific bibliographic reference for that research.
\end{itemize}
\end{document}
